%!TEX Root = ../main.tex
\chapter{Associated Primes and Primary Decomposition}
In this chapter, we will discuss primary decomposition,
which is a theory that tries to decompose a general ring
into a simpler set of ideals, called \vocab{primary ideals}.

Many theorems in this chapter requires Noetherian assumption.
We will try to be careful about which part requires
Noetherian hypothesis, but for most part,
you should assume that all rings are Noetherian.

\section{Associated Primes}
\subsection{Definition}
\begin{definition}[Associated primes]
    Let $M$ be an $A$-module.
    A prime ideal $\pp\subset A$ is an \vocab{associated prime} $M$
    if there exists $x\in M$ for which
    $\Ann(x) = \pp$.
    \notation{$\Ass(M)$}{associated primes}

    The set of associated primes is (unfortunately) denoted $\Ass(M)$
    (or $\Ass_A(M)$ if the ground ring is not clear).
\end{definition}
Equivalent definitions: $\pp\in\Ass(M)$ iff
the following (equivalent) conditions hold:
\begin{itemize}
    \item there exists $x\in M$ such that the map 
    $A\stackrel{\cdot x}\longto M$ has kernel $\pp$.
    \item there exists an injective map $A/\pp\to M$.
\end{itemize}

Let's do some examples.
\begin{example}
    We have some examples for $\mathbb Z$ module.
    \begin{itemize}
    \item If $M=\mathbb Z/(p^aq^b)$, then
	$\Ass(M) = \{(p), (q)\}$,
	where $(p) = \Ann(p^{a-1}q^b)$, $(q) = \Ann(p^aq^{b-1})$.
    \item If $M=\mathbb Z\oplus\mathbb Z/2$,
	then $\Ass(M) = \{(0), (2)\}$.
    \end{itemize}
\end{example}
\begin{exercise}
Let $M$ and $N$ be two $R$-modules.
Prove that $\Ass(M\oplus N) = \Ass(M)\cup \Ass(N)$.
\end{exercise}
%
%
%
%
\subsection{Basic Properties}
\begin{proposition}
    \label{prop:basic_ass}
    We have the following.
    \begin{enumerate}[label=(\alph*)]
    \item If $x\in M$ and $\Ann(x)=\pp$.
	Suppose that $0\neq y\in Ax$, then $\Ann(y)=\pp$ also.
    \item $\Ass(A/\pp) = \{\pp\}$.
   \end{enumerate}
\end{proposition}
\begin{proof}
    \begin{enumerate}[label=(\alph*)]
    \item We have $Ax \simeq A/\pp\subset M$,
	and since $A/\pp$ is a domain and $y\neq 0$,
	we have $\Ann_{A/\pp} y = (0)$.
	Thus, $\Ann_A y = \pp$.
    \item Follows directly from (a). \qedhere
    \end{enumerate}
\end{proof}

We now prove that there exists at least one associated prime.
\begin{theorem}[Existence]
    \label{thm:ass_exists}
    If $A$ is Noetherian and $y\neq 0\in M$, 
    then there exists a prime $\pp\in \Ass(M)$
    such that $\pp\supseteq \Ann(y)$.

    \noindent 
    In particular, if $M\neq 0$, then $\Ass(M)\neq\emptyset$.
\end{theorem}
\begin{proof}  
Make a poset
$$\Sigma = \{\Ann(x) : x\neq 0\in M\}$$
ordered by inclusion.

We claim that every maximal element must be in $\Ass(M)$.
To show this, suppose $I\in\Sigma$ is maximal element,
so $I=\Ann(x)$ for some $x\in M$. 
It suffices to that $I$ is prime.
To that end, given $f,g\in A$ and $fg\in I$.
Then, $(fg)x=0$. 

Assume that $g\notin I$. Thus, $gx\neq 0$.
Then, we have that
$(I,f)\subset \Ann(gx) \in\Sigma,$
but $I$ is maximal in $\Sigma$, 
so the only way this could happen is $\Ann(gx)=I$ and $f\in I$.

Finally, if $\Ann(y)$ is not maximal in $\Sigma$,
we keep climbing up the poset, forming an infinite
ascending chain that must stabilize because $A$ is Noetherian.
The poset must stabilize at a maximal element of $\Sigma$,
which is in $\Ass(M)$ by the above claim.
\end{proof}
\begin{remark}
    \Cref{thm:ass_exists} does not hold when $A$ is not Noetherian.
    Example: $R=M=C(\RR,\RR)$, the ring of continuous functions
    on a real line.
\end{remark}
\begin{corollary}
    \label{cor:union_ass}
    We have
    $$\bigcup_{\pp\in \Ass(M)} \pp = \bigcup_{y\neq 0\in M}
     \Ann(y) = \{\text{zero divisors of } M\}.$$
\end{corollary}
We now show that for a Noetherian module,
there are finitely many associated primes.
\begin{theorem}[Finiteness]
\label{thm:ass_is_finite}
\begin{alphenum}
    \item If we have a short exact sequence of $A$-modules
    $$0\longto L\stackrel\phi\longto M \stackrel\pi \longto N\longto 0,$$ 
    then
    $$\Ass(L)\ \subseteq\ \Ass(M)\ \subseteq\ \Ass(L)\cup \Ass(N).$$
    \item If $M$ is a Noetherian $A$-module,
    then $\Ass(M)$ is finite.
\end{alphenum}
\end{theorem}
\begin{proof}
\begin{alphenum}
\item 
For the first containment, we have
the composition $0\to A/\pp\to L\stackrel\phi\to M$ of injective maps,
therefore, $\pp\in\Ass(M)$ also.

We now prove the second containment.
Let $\pp\in\Ass(M)$, so there exists injective map
$\alpha : A/\pp \to M$. By viewing $L$ as a submodule of $M$,
we assume that $L=\Ker\pi$.  We have two cases.
\begin{itemize}
\item \textbf{If $\Im(\alpha)\cap L = 0$},
then $\pi\circ\alpha$ is an injective from
$A/\pp\to N$, and so $\pp\in\Ass(N)$.
\item \textbf{Otherwise}, there exists $y\in\Im(\alpha)\cap L$.
By \Cref{prop:basic_ass} (a), we have that $\Ann(y) = \pp$, so $\pp\in\Ass(L)$.
\end{itemize}
\item 
It suffices to make a chain of submodules
$$0=M_0\subseteq M_1 \subseteq M_2 \subseteq \dots \subseteq M_n
= M$$
such that $M_i/M_{i-1} = A/\pp_i$ for prime ideal $\pp_i$.
If we have this, then (a) gives $\Ass(M)\subseteq \{\pp_1,\dots,\pp_n\}$.

To prove this, by \Cref{thm:ass_exists},
pick $\pp_1\in\Ass(M)$, so
there exists element $x_1\in M$ such that $\Ann_M(x_1)=\pp_1$.
Thus, we let $M_1 = Ax_1 \simeq A/\pp_1$.

If $M_1=M$, then we are done.
Otherwise, by \Cref{thm:ass_exists}, pick $\pp_2 \in \Ass(M/M_1)$,
so there exists element $x_2\in M$ such that 
$\Ann_{M/M_1}(x_2) = \pp_2$. Therefore, $M_2 = Ax_2$
satisfies $M_2/M_1 = A/\pp_2$. 
Continue this process forever, and it must terminate because
$M$ is a Noetherian $A$-module. \qedhere
\end{alphenum}
\end{proof}
\begin{example}
    Consider the exact sequence
    $$0\to\mathbb Z\stackrel{\cdot 2}{\longrightarrow}\mathbb Z
    \to \mathbb Z/2\to 0.$$
    Then, we have $\Ass M = \{(0)\} \neq \Ass(L)\cup\Ass(N)
    =  \{(0), (2)\}$.
\end{example}
\begin{theorem}
    \label{thm:ann_ass}
    \begin{alphenum}
    \item For any multiplicative set $S\subseteq A$, we have
    $$\Ass_{S^{-1}A}(S^{-1}M) = \{S^{-1}\pp : 
    \pp\in \Ass_A(M)\text{ such that }\pp\cap S=\emptyset\}.$$
    \item If $A$ is Noetherian, $M$ is finitely-generated, and $\pp$ is a prime ideal 
    minimal among those containing $\Ann(M)$,
    then $\pp\in\Ass(M)$.
    \end{alphenum}
\end{theorem}
\begin{proof}
    \begin{alphenum}
    \item Mostly definition-chasing.
    
    If $\pp\in\Ass_A(M)$ and $\pp\cap S = \emptyset$,
    then there is an injective map $R/\pp \hookrightarrow M$,
    which turns to $S^{-1}(A/\pp) = S^{-1}A/S^{-1}\pp 
    \hookrightarrow S^{-1}M$,
    where we used \Cref{thm:localization_exact}
    to commute localization inside quotient and 
    see that it preserves injectivity.
    This gives $S^{-1}\pp \in \Ass_{S^{-1}A}(S^{-1}M)$.

    Conversely, by \Cref{thm:prime_ideal_localize},
    suppose that $S^{-1}\pp \in \Ass_{S^{-1}A}(S^{-1}M)$,
    where $\pp\cap S=\emptyset$,
    then we have an injective map 
    $S^{-1}(A/\pp) \hookrightarrow S^{-1}M$.
    Suppose that this map takes $1$ to $x/s$ for some $x\in M$
    and $s\in S$. 
    We claim that $\Ann(x) = \pp$.
    To see this, suppose that $rx=0$ for some $r\in A$.
    Then, $r\tfrac xs = 0$, so $r=0$ in $S^{-1}(A/\pp)$.
    Therefore, $ur = 0\pmod\pp$ for some $u\in S$,
    and since $S\cap\pp\neq\emptyset$,
    we have $u\notin\pp$, so $r\in\pp$.

    \item First, we show that $M_\pp\neq 0$.
    Suppose $M_\pp=0$. Suppose $x_1,x_2,\dots,x_n$ generates $M$.
    Then, since the image of $x_i$ in $M_\pp$ is $0$,
    there exists $a_i\notin\pp$ such that $a_ix_i=0$.
    Now, note that the product $a_1\cdots a_n$ annihilates $x_1,\dots,x_n$,
    and hence annihilates $M$, so $a_1\cdots a_n \in \Ann(M)$,
    but $a_1\cdots a_n\notin\pp$, a contradiction.
    
    By \Cref{thm:ass_exists}, we have 
    $\Ass_{A_\pp}(M_\pp) \neq \emptyset$.
    However, from (a), this set correspond to $\qq\in \Ass(M)$
    such that $\qq\subseteq\pp$. Note that $\qq$ contains $\Ann(M)$,
    so by minimality, $\qq=\pp$. Hence $\pp\in\Ass(M)$.
    \qedhere
    \end{alphenum}
\end{proof}
%
%
%
%
\subsection{Local Property and Associated Primes}
\begin{corollary}[Local Property and Associated Primes]
    \label{cor:local_element_ass}
    Let $A$ be Noetherian ring, $M$ be an $A$-module.
    Suppose that $x\in M$ such that for all $\pp\in\Ass(M)$,
    the image of $x$ in $M_\pp$ is $0$.
    Then, $x=0$.
\end{corollary}
\begin{proof}
    Assume $x\neq 0$.
    By \Cref{thm:ass_exists}, $\Ann(x)$ is contained in
    some associated prime $\pp$.
    Then, if the image of $x$ in $M_\pp$ is $0$,
    then $ax=0$ for some $a\notin\pp$.
    This contradicts $\Ann(x) \subseteq\pp$.
\end{proof}
This gives the following two corollaries.
\begin{corollary}  
    \label{cor:local_ass}
    Let $A$ be Noetherian and $M$ be an $A$-module.
    Then,
    \begin{alphenum}
    \item If $K$ is a submodule of $M$, 
        then $K=0$ if and only if $K_\pp=0$
	for all $\pp\in \Ass(M)$.
    \item If $\phi : M\to N$, then 
	$\phi$ is injective
	if and only if $\phi_\pp : M_\pp \to N_\pp$
	is injective for $\pp\in\Ass(M)$.
    \end{alphenum}
\end{corollary}
\begin{proof}
    \begin{alphenum}
    \item Clear from \Cref{cor:local_element_ass}.
    \item Apply (a) with $K=\Ker\phi\subseteq M$. \qedhere
    \end{alphenum}
\end{proof}
%
%
%
%
%
%
\section{Primary Decomposition}
\label{sec:primary_decomp}
\begin{center}
    \fbox{For this entire section, let $A$ be a
    a Noetherian ring.}
\end{center}
The goal of this section is to develop 
the theory that decompose ideals of $A$ into simpler ideals.
We have seen two special cases of this.
\begin{itemize}
    \item If $A$ is a PID (hence UFD), then 
    an ideal $I=(f)$ can be written in form
    $$I = (p_1)^{a_1} \cdots (p_k)^{a_k}
    = \bigcap_{i=1}^k (p_i)^{a_i},$$
    where $p_1,\dots,p_k$ are prime elements of $A$.
    In order words, any ideal is an intersection of 
    prime powers. 
    Special cases of this include $A=\ZZ$
    and $A=k[x]$.
    \item From \Cref{subsec:irreducible_comp},
    any radical ideal $I$ is an intersection of 
    prime ideals minimal among those containing $I$.
\end{itemize}
Our theorem will be that for any Noetherian ring $A$,
any ideal $I$ can be written as the intersection
of a special kind of ideas called primary ideals.
When $A$ is a PID, primary ideals coincide with 
prime powers. 
However, the notion of primary ideals is more general.
%
%
%
\subsection{Primary Ideals}
\begin{definition}
    An ideal $Q\subsetneq A$ is 
    \begin{itemize}
        \item \vocab[primary ideals]{primary}
    if whenever $fg\in Q$, either $f\in Q$ or $g^n\in Q$
    for some $n\geq 1$.
        \item \vocab{$\pp$-primary} (where $\pp$ is a prime ideal) if $Q$ is primary and $\sqrt Q=\pp$.
    \end{itemize}
\end{definition}
In other words, $Q$ is primary if and only if
every zero divisor in $A/Q$ is nilpotent.
For example, prime powers in $\mathbb Z$ are primary.
\begin{lemma}
    If $Q$ is primary, then $\sqrt Q$ is prime.
\end{lemma}
\begin{proof}
    Note that 
    $$fg\in\sqrt Q\implies f^ng^n\in Q
    \implies f^n\in Q\text{ or } g^{mn}\in Q.$$
    Either way, this implies that $f\in\sqrt Q$
    or $g\in\sqrt Q$.
\end{proof}
In the definition above, since we assumed that $A$ is Noetherian,
$\pp$ is finitely generated,
so there exists $m>0$ such that $\pp^m\subset Q\subset\pp$.
\begin{warn}
\label{warn:primary_ideals}
\begin{itemize}
\item \textbf{If $\sqrt Q$ is prime, then $Q$ is 
not necessarily primary.}
For example,
$$I = (x^2,xy)\subset\mathbb C[x,y] \implies \sqrt I = (x).$$
We have $(x^2)\subset I\subset (x)$, but $I$ is not primary.
We have $xy\in I$ but $x\notin I$ and $y\notin\sqrt I$.
\item \textbf{Prime powers are not necessarily
primary}. To see an example, take
$$A = \frac{\mathbb C[x,y,z]}{(xz-y^2)},\ \pp=(x,y).$$
Then, we claim that $\pp^2 = (x^2, xy, y^2)$ is not primary.
To see why, note that $xz = y^2$,
but $x\notin\pp^2$ and $z\notin\pp$.
\end{itemize}
\end{warn}
\begin{proposition}
    \label{prop:max_ideal_primary}
    Given an ideal $Q\subset A$ such that $\sqrt Q = \mathfrak m$
    is maximal, then $Q$ is $\mathfrak m$-primary.
\end{proposition}
\begin{proof}
    Suppose that $f\notin Q$, then the ideal
    $I =\{g\in A : fg\in Q\}$ is a proper ideal of $A$,
    so it is contained in a maximal ideal $\mathfrak m'$.
    In particular, we have $Q\subset I\subset\mathfrak m'$.
    Thus, we have $\sqrt Q\subset\sqrt{\mathfrak m'}
    \implies \mathfrak m \subset\mathfrak m'$,
    and so $\mathfrak m = \mathfrak m'$.
    Thus, $I\subset\mathfrak m$.
    Therefore, if $fg\in Q$, we must have $g\in I$,
    so $g\in\mathfrak m$
    or $g^n\in Q$ for some $I$.
\end{proof}
We now connect primary ideal to associated primes
we studied earlier in the previous section.
\begin{theorem}[Primary if and only if one associated prime]
    \label{thm:primary_one_ass}
    $Q$ is $\pp$-primary if and only if 
    $\Ass(A/Q) = \{\mathfrak p\}$.
\end{theorem}
\begin{proof}
    $(\Rightarrow)$. 
    Pick $x\in A$ such that
    the image $\overline x\in A/Q$ is nonzero and
    such that $\Ann_{A/Q}(\overline x)$ is a prime ideal.
    Then, we observe that
    \begin{align*}
        Q \subseteq \Ann_{A/Q}(\overline x) 
        &= \{y \in A : xy\in Q\} \\
        &\subseteq \{y : y^n\in Q\text{ for some }n\geq 0\} 
        \tag{def. of primary ideal}\\
        &\subseteq \sqrt{Q} = \pp.
    \end{align*}
    Since $\Ann_{A/Q}(\overline x)$ is prime ideal,
    taking radical both sides gives 
    $\pp\subseteq \Ann_{A/Q}(\overline x)
    \subseteq \pp$, so $\Ann_{A/Q}(\overline x) = \pp$,
    so any associated prime must be $\pp$.
    Since $\Ass(A/Q)\neq\emptyset$ (\Cref{thm:ass_exists}),
    it follows that $\Ass(A/Q)=\{\pp\}$.
    
    ($\Leftarrow$) If $\Ass(A/Q) = \{\pp\}$,
    then any nonzero submodule of $A/Q$ also has only $\pp$
    as associated prime.
    Given $0\neq M\subset A/Q$, we note by definition 
    of radical that
    \begin{equation}
    \label{eq:primary_ass}
    \sqrt{\Ann M} = \bigcap_{\qq\supseteq\Ann(M)} \qq
    = \bigcap_{\qq\in\Ass(M)}\qq = \pp
    \end{equation}
    (where we used \Cref{thm:ann_ass} in the middle equality).
    In particular, plugging $M=A/Q$ into \eqref{eq:primary_ass} gives $\sqrt Q = \pp$.

    Now, given $f,g$ such that $fg\in Q$.
    If $f\notin Q$, then we reduce
    $0\neq \overline f\in A/Q$,
    giving a submodule $M=A\overline f \subset A/Q$ 
    generated by single element $\overline f$
    We have $g\in\Ann(M)$, so using \eqref{eq:primary_ass}
    with this choice of $M$, we have
    $g\in\sqrt{\Ann(M)} = \pp$, so $g^n\in\pp$
    for some $n$. Hence, $Q$ is primary.
\end{proof}
\begin{remark}
    All the theory above generalizes by replacing
    $A$ with arbitrary modules $M$.
    In that case, a submodule $N\subseteq M$
    is $\pp$-primary if $\Ass(M/N) = \{\pp\}$.
    The theorem then says that any submodule of $M$
    can be expressed as an intersection of $\pp$-primary
    submodules.
    Moreover, the set of primes that appear is precisely
    $\Ass(M)$.
\end{remark}
%
%
%
%
%
\subsection{Primary Decomposition}
We are now ready to state primary decomposition theorem.
\begin{definition}[Primary decoposition]
    Given $I\subset A$, a \vocab{primary decomposition}
    of $I$ is an expression
    $$I=\bigcap_{i=1}^n Q_i$$
    such that
    \begin{itemize}
    \item $Q_i$ is $\pp_i$-primary for all $i$.
    \item 
    $I\not\subset\bigcap_{i\neq j} Q_i$ for each $j$.
    \item $\pp_i\neq\pp_j$ for all $i,j$.
    \end{itemize}
\end{definition}
The second one is to clean out redundant elements.
For the third one, if $\pp_i=\pp_j$ for some $i$ and $j$,
then the following lemma tells us that 
we can combine $Q_i$ and $Q_j$ into a single ideal.
\begin{lemma}
    If $Q_1$ and $Q_2$ are both $\pp$-primary,
    then so is $Q_1\cap Q_2$.
\end{lemma}
\begin{proof}
    The map $A/(Q_1\cap Q_2) \longto A/Q_1\oplus A/Q_2$
    is injective. Thus, by \Cref{thm:ass_is_finite} (a),
    $$\Ass(A/(Q_1\cap Q_2))\ \subseteq\ \Ass(A/Q_1)
    \cup \Ass(A/Q_2) \ =\ \{\pp\},$$
    so since $\Ass(A/(Q_1\cap Q_2))\neq\emptyset$
    (\Cref{thm:ass_exists}),
    it follows that $\Ass(A/(Q_1\cap Q_2)) = \{\pp\}$,
    so the result follows from \Cref{thm:primary_one_ass}.
\end{proof}
Thus, to prove that primary decomposition exists,
we must show that any ideal can be written as
intersection of primary ideals; there is no need to care
about the second and third condition.

Let us demonstrate some examples of what primary 
decomposition look like.
\begin{example}
Consider
$$I = (x^2,xy)\subset\mathbb C[x,y].$$
This ideal is contained in the maximal ideal $(x,y)$, and in fact, we have
$$(x^2,xy) = (x)\cap (x,y)^2,$$
which is one primary decomposition.
However, there are also other decompositions. For example,
\begin{align*}
    I &= (x)\cap (x^2,y) \\
      &= (x)\cap (x^2 y-\alpha x).
\end{align*}
\end{example}
\begin{example}
    Consider
    $$A = \frac{\mathbb C[x,y,z]}{(xz-y^2)},\quad \pp=(x,y).$$
    We have seen in \Cref{warn:primary_ideals} that
    $\pp^2$ is not primary.
    Its primary decomposition looks like
    $$\pp^2 = (x)\cap (x,y,z)^2.$$
\end{example}
We now prove that primary decomposition exists.
\begin{theorem}
    If $A$ is Noetherian, then primary decomposition always exists.
\end{theorem}
To prove this theorem, we first make the definition.
\begin{definition}[Indecomposable ideals]
    An ideal $I\subset A$ is \vocab{indecomposable}
    if whenever $I = J\cap K$ for ideals $J$ and $K$,
    then $I=J$ or $I=K$.
\end{definition}
Thus, by repeatedly breaking down decomposable ideals,
any ideal $I$ must be written as $J_1\cap J_2\cap\dots\cap J_n$,
where $J_i$ indecomposable.
Since $A$ is Noetherian, the process must stop.

Thus, to prove that primary decomposition exists,
it suffices to show the following.
\begin{proposition}[Indecomposable implies primary]
    If $Q$ is indecomposable ideal
    of a Noetherian ring $A$, then $Q$ is primary.
\end{proposition}
\begin{proof}
    By replacing $A$ with $A/Q$, we may assume that $Q=(0)$.

    If $xy=0$, then we have
    $$y\in\Ann(x)\subset\Ann(x^2)\subset \Ann(x^3)\subset\dots$$
    This chain must stabilize, so for some $k$, we have
    $\Ann(x^k) = \Ann(x^{k+1})$.

    Now, we claim $(x^k)\cap (y)=(0)$.
    To see why, if $a\in (x^k)\cap (y)$, then
    $ax=0$ and $a=bx^k$. Thus, $bx^{k+1}=0\implies
    b\in\Ann(x^{k+1}) = \Ann(x^k)$.
    Thus, $a=bx^k = 0$ as claimed.

    Since $(0)$ is indecomposable, we must have
    either $(x^k) = (0)$ or $(y)=(0)$.
    Thus,  $x^k=0$ or $y=0$, as desired.
\end{proof}
%
%
%
%
%
\subsection{Quasi-uniqueness of Primary Decomposition}
Unfortunately, as seen above,
primary decompositions are not unique.
However, we can say something about uniqueness.
\begin{theorem}[Quasi-uniqueness]
    Let $A$ be a Noetherian ring.
    If an ideal $I$ has a primary decomposition $I=\bigcap Q_i$
    such that $Q_i$ is $\pp_i$-primary,
    then $\{\pp_1,\dots,\pp_n\} = \Ass(A/I)$.

    In particular, the set $\{\pp_1,\dots,\pp_n\}$
    is unique.
\end{theorem}
\begin{proof}
    Consider the direct sum of projection
    $$f : A\to \bigoplus_{i=1}^n A/Q_i.$$
    We have $\ker f = I$. Thus, $A/I$
    is a subring of these direct sums.
    In particular,
    $$\Ass(A/I) \subseteq \Ass\left(\bigoplus_{i=1}^n A/Q_i\right)
    = \bigcup_{i=1}^n \Ass(A/Q_i) = \{\pp_1,\dots,\pp_n\}.$$

    We now show that for each $i$, we have $\pp_i\in \Ass(A/I)$.
    Let
    $$N = \frac{\bigcap_{j\neq i} Q_j}{I}\subset A/I.$$
    By irredundant condition, 
    $N\neq 0$. Moreover, we have the composition
    $N\to A/I\to \bigoplus_{j=1}^n A/Q_j$ is zero on each factor $j\neq i$,
    and it is injective onto $A/Q_i$.
    Thus, $\Ass(N) \subset \Ass(A/Q_i) =  \{\pp_i\}$,
    so $\Ass(N) = \{\pp_i\}$ (because it must not be empty),
    and so $\pp_i\in\Ass(A/I)$.
\end{proof}
\section{Problems}
\begin{problem}
Let $R=\mathbb C[x,y]/(xy)$-module and let $M$ be an $R$-module with generators $a$ and $b$ and a single relation $xa+yb=0$ (the same module as in \Cref{prob:diff_xy}). Determine $\operatorname{Ass} M$.
\end{problem}
\begin{problem}
Let $A$ be a Noetherian ring and $M$ be a finitely-generated module. Suppose that $I$ is an ideal of $A$ with property that for every $a\in I$, there exists $0\neq m_a\in M$ such that $am_a = 0$. Prove that there exists $0\neq m\in M$ such that $I\cdot m=0$.

Hint: prime avoidance, \Cref{prob:prime_avoidance}.
\end{problem}
\begin{problem} 
\label{prob:ass_reduced}
Let $R$ be a Noetherian ring and $I$ be an ideal.
\begin{alphenum}
\item Explain why every prime ideal minimal over $I$ 
(i.e., minimal among those containing $I$)
 must be an associated prime of $R/I$.
\item If $R/I$ is reduced 
(i.e., $I$ is a radical ideal), 
prove that every associated prime of $R/I$ is minimal over $I$. 
\end{alphenum}
\end{problem}
\begin{problem}
\label{prob:ass_denominator}
Let $A$ be a Noetherian domain and $r\in\Frac A$.
Suppose that for every $f\in A$ and $\pp\in \Ass(A/(f))$,
we have $r\in A_\pp$. Prove that $r\in A$.

(Hint: suppose that $r=a/b$. Let $M$ be a submodule of 
$A/(b)$ generated by $a$ and consider $\pp\in \Ass M$.)
\end{problem}
\begin{problem}
\label{prob:primary_quotient_localization}
Let $R$ be a ring and $Q$ be a primary ideal.
\begin{alphenum}
\item If $\phi : S\to R$ is a ring homomorphism,
then prove that $\phi^{-1}(Q)$ is a primary ideal in $S$.
\item Prove that for every multiplicative set 
$S\subseteq R$, $S^{-1}Q$ is either equal to $S^{-1}R$
or a primary ideal in $S^{-1}R$.
\end{alphenum}
\end{problem}
\begin{problem}
Find a primary decomposition of the following ideals.
\begin{alphenum}
\item Ideal $(x+y,xy)$ of $\mathbb C[x,y]$.
\item Ideal $(xy, x-yz)$ of $\mathbb C[x,y,z]$.
\item Ideal $(x^2,xy^2,y^2z,xz)$ of $\mathbb C[x,y,z]$.
\end{alphenum}
(\Cref{prob:primary_quotient_localization} might be helpful.)
\end{problem}